% Terjemahan dibuat oleh AI. Koreksi dan tata letak dihasilkan oleh OpenAI Codex — GPT-6 Astra, tingkat penalaran Ultra. Nama Emmy Noether dan para penulis yang tercantum pada karya tetap dipertahankan. Tidak ada penerjemahan manusia atau penelaahan ilmiah manusia independen yang diatribusikan kepada penerbit.
% Hak tetap dimiliki pemegangnya masing-masing. Penyebaran gratis untuk tujuan nonkomersial dan murni pendidikan didasarkan pada anggapan penerbit bahwa pemegang hak tidak keberatan, bukan izin yang telah dipastikan. Akses gratis tidak dengan sendirinya mengizinkan publikasi turunan. Pemegang hak dapat mengajukan keberatan melalui tautan berikut; penerbit akan menarik materi terkait atas permintaan mereka. Tidak ada lisensi baru atau penyerahan ke domain publik.
% https://github.com/KokunoYumeto/modern-latex-manuscripts/issues/new
\documentclass[11pt]{article}
\usepackage[a4paper,margin=2.4cm]{geometry}
\usepackage{fontspec}
\setmainfont{Noto Serif}
\setlength{\parindent}{1.2em}
\setlength{\parskip}{0.5em}
\setlength{\emergencystretch}{2em}
\begin{document}

\section*{36. Pendiferensialan Ideal dan \emph{Differente}}

\begin{center}
\emph{J. Ber. d. DMV 39 (1930), hlm. 17}
\end{center}

2. E. Noether, Göttingen: Pendiferensialan ideal dan \emph{Differente}.

Ditunjukkan bagaimana \emph{Differente} dari suatu medan bilangan aljabar dapat dipandang sebagai \emph{hasil bagi diferensial} dari suatu \emph{ideal} yang mendefinisikan medan tersebut. Bahwa definisi ini bertepatan dengan definisi yang lazim merupakan konsekuensi dari teorema-teorema struktur yang juga menarik secara tersendiri; dalam arti yang analog, teorema-teorema itu merupakan perampatan rumus interpolasi Lagrange.

Uraian terperinci akan dimuat dalam Math. Ann.

\end{document}
